% =====================================================================
%  c03 —— 植田一石《同调镜像对称》中译（第 6 节末 + 第 7–9 节开篇）
%  页范围：p014–p015
%  编号公式：9 条 —— (6.3) (6.4) (7.1)–(7.5) (8.1) (9.1)
%  说明：BRIEF 原记 10 条，系 work/tags.txt 把 p014 正文中的交叉引用
%        「(6.3) において」误计为编号公式；按页面图像实际为 9 条。
%  节标题：第 7 节（起于 p014）、第 8 节（起于 p015）、第 9 节（起于 p015）
%        均在本 chunk 内，故各写一次；p014 顶部为第 6 节续文，不再写节标题。
% =====================================================================

粘合而得到的流形。$\mathrm{Sp}(2n,\R)$ 结构的可积性与 $\omega$ 为闭形式等价（Darboux 定理），$\mathrm{GL}(n,\C)$ 结构的可积性与 Nijenhuis 张量消失等价（Newlander–Nirenberg 定理）。

复几何与辛几何乍看之下是完全不同类型的几何，但在某种意义上非常相似\footnote{symplectic 一词，是 Weyl 把源自拉丁语的 complex 一词的构成要素 ‘com’ 与 ‘plexus’ 替换为对应的希腊语 ‘sym’ 与 ‘plectikos’ 而创造出来的。}。二者都是偶数维的，可以认为它们在某种意义上是把实流形「加倍」而成\footnote{此外，作为统一处理这两种几何的框架，Hitchin 引入了广义复结构（generalized complex structure）的概念。关于此可参见例如 [Hit11]。}。复流形的典型例子是切丛，辛流形的典型例子是余切丛。镜像对称给出这两种几何之间的关系。

\begin{equation}
    \mathrm{GL}(n,\C)\cap\mathrm{Sp}(2n,\R)\cap O(2n)=U(n) \tag{6.3}
\end{equation}

因此，$U(n)$ 结构诱导出 $\mathrm{GL}(n,\C)$ 结构与 $\mathrm{Sp}(2n,\R)$ 结构，而当这个 $\mathrm{GL}(n,\C)$ 结构与 $\mathrm{Sp}(2n,\R)$ 结构都可积时，称这个 $U(n)$ 结构为 Kähler 结构（Kähler structure）。在 (6.3) 中，任意两个的公共部分自动包含在剩下的一个之中，这对应于：在 Kähler 流形中，复结构、辛结构、Riemann 结构中的任意两个都通过
\begin{equation}
    g(V,W)=\omega(V,JW) \tag{6.4}
\end{equation}
来确定剩下的一个。

\section{分析力学}

分析力学（analytical mechanics）大致分为 Lagrange 形式（Lagrangian formalism）与 Hamilton 形式（Hamiltonian formalism）两种。在 Lagrange 形式的分析力学中，系统由称为构型空间（configuration space）的流形 $N$、以及其切丛 $TN$ 上称为 Lagrange 函数（Lagrangian）的函数 $L(q,v)$ 指定。这里 $q$ 是 $N$ 的局部坐标，$v$ 是对应于 $q$ 的 $TN$ 的纤维方向坐标。系统的时间演化被定义为 $N$ 的路径空间（path space）$\cP N:=\{\gamma\colon[0,1]\to N\}$ 上的作用泛函（action functional）
\begin{equation}
    S\colon\cP N\to\R,\qquad \gamma\mapsto\int L(\gamma,\dot\gamma)\,dt \tag{7.1}
\end{equation}
的驻点，并由 Euler–Lagrange 方程
\begin{equation}
    \frac{d}{dt}\left(\left.\frac{\partial L(q,v)}{\partial v}\right|_{v=\dot q}\right)-\left.\frac{\partial L(q,v)}{\partial q}\right|_{v=\dot q}=0 \tag{7.2}
\end{equation}
刻画。% [Tutor Note: 原文 p014 此句作「Hailton 形式の解析力学」，Hailton 当为 Hamilton 之笔误；译文按 Hamilton 处理，请审校确认。]
另一方面，在 Hamilton 形式的分析力学中，系统由辛流形 $(M,\omega)$ 与 $M$ 上称为 Hamilton 函数（Hamiltonian）的函数 $H(p,q)$ 指定。这里 $(p,q)$ 是 $M$ 的 Darboux 坐标。系统的时间演化被定义为 $H$ 生成的 Hamilton 向量场（Hamiltonian vector field）
\begin{equation}
    X_H:=\sum_{i=1}^{n}\left(\frac{\partial H}{\partial p_i}\frac{\partial}{\partial q_i}-\frac{\partial H}{\partial q_i}\frac{\partial}{\partial p_i}\right) \tag{7.3}
\end{equation}
的积分曲线（integral curve），并被刻画为作用泛函
\begin{equation}
    S=\int (p\,dq-H\,dt) \tag{7.4}
\end{equation}
的驻点。

将 Lagrange 函数作 Legendre 变换而得到的 $T^*N$ 上的函数
\begin{equation}
    p:=\frac{\partial L(q,v)}{\partial v},\qquad H(q,p):=p\cdot v-L \tag{7.5}
\end{equation}
取作 Hamilton 函数，则可以把 Lagrange 形式的分析力学改写为 Hamilton 形式\footnote{辛流形不总是余切丛，因此 Hamilton 形式比 Lagrange 形式自由度更大。另一方面，若辛流形 $(M,\omega)$ 具有拉格朗日子流形 $L$，则 $M$ 中 $L$ 的管状邻域与 $L$ 的余切丛 $T^*L$ 辛同胚（Weinstein 邻域定理（Weinstein neighborhood theorem））。}。此事虽与上述并不直接相关，但 Legendre 变换在镜像对称中起着重要作用。

\section{Liouville–Arnold 定理}

Hamilton 系统若具有守恒量，则通过将其取为 Darboux 坐标的一部分，系统的维数实质上降低 2。当守恒量有多个时，要能把它们同时取为 Darboux 坐标的一部分，它们必须相互 Poisson 交换且函数独立。这里，函数 $f_1,\dots,f_k$ 称为函数独立（functionally independent），是指 $df_1,\dots,df_k$ 在某个稠密子集上线性无关。Poisson 交换且函数独立的守恒量的最大个数不超过维数的一半，而这种个数恰好为维数一半的 Hamilton 系统称为完全可积系统。

设 $\pi:=(f_1,\dots,f_n)\colon X\to\R^n$ 为完全可积系统的守恒量所定的映射，$b\in\R^n$ 为其正则值。$\pi$ 的纤维上，由这些守恒量所伴随的 Hamilton 流自然带入 $\R^n$ 的作用，关于这个作用成为齐性空间。$\R^n$ 的紧连通齐性空间只限于环面，因此若 $\pi$ 的纤维紧且连通，就是环面。令 $\pi$ 的像为 $B$，正则值的集合为 $B^{\mathrm{sm}}\subset B$。若 $\pi$ 是固有映射，且所有纤维都连通，则由 Ehresmann 纤维化定理，$\pi|_{\pi^{-1}(B^{\mathrm{sm}})}\colon\pi^{-1}(B^{\mathrm{sm}})\to B^{\mathrm{sm}}$ 成为以环面为纤维的纤维丛。这个纤维丛的纤维的切空间由守恒量生成的 Hamilton 向量场张成，而守恒量的个数是全空间维数的一半，且这些 Hamilton 向量场相互交换，所以纤维成为拉格朗日子流形。取 $B^{\mathrm{sm}}$ 的充分小的开集 $U$ 与 $U$ 上 $\pi$ 的平凡化，把 $\pi|_{\pi^{-1}(U)}\colon\pi^{-1}(U)\to U$ 与投影 $\pi_1\colon U\times L\to U$ 等同。任意取定 $H_1(L;\Z)$ 的 $\Z$ 上有序基 $(\alpha_i)_{i=1}^{n}$、各 $\alpha_i$ 所代表的 $L$ 的循环 $C_i$、以及 $U$ 的基点 $b_0$。对任意 $b\in U$，任意取 $U$ 上从 $b_0$ 到 $b$ 的道路 $\gamma\colon[0,1]\to U$，$\gamma(0)=b_0$，$\gamma(1)=b$，则 $\gamma([0,1])\times C_i$ 定出 $U\times L$ 的一个 2-cycle。这个 2-cycle 上的积分
\begin{equation}
    x_i(b):=\int_{\gamma([0,1])\times C_i}\omega \tag{8.1}
\end{equation}
只依赖于 $b$ 与基 $(\alpha_i)_{i=1}^{n}$（由于 $\omega$ 是闭形式，故不依赖于 $\gamma$ 的取法；由于 $\omega|_L$ 为 0，故不依赖于 $\alpha_i$ 的代表元 $C_i$ 的取法）。这样定出的函数 $(x_i)_{i=1}^{n}\colon U\to\R^n$ 给出 $U$ 的坐标，称为 $U$ 的作用坐标（action coordinate）。伴随这个作用坐标的 $\R^n$ 的 Hamilton 作用，其固定子群为 $\R^n$ 的自然格 $\Z^n$，因此任意取定 $L$ 的基点 $l_0$ 后，就引起各纤维 $\pi^{-1}(b)$ 与 $\R^n/\Z^n$ 的等同。这样得到的纤维的坐标称为角坐标（angle coordinate）。进一步，系统原来的 Hamilton 函数只是作用坐标的函数，并且可知：

\begin{itemize}[leftmargin=2.4em]
    \item 正则水平集紧致且为连通分支时，它是环面；
    \item 这个环面上带有称为作用–角坐标的特殊坐标；
    \item 系统的时间演化关于作用坐标是常数，关于角坐标是线性的，
\end{itemize}

这称为 Liouville–Arnold 定理。即，完全可积系统是与拉格朗日环面纤维丛几乎等价的概念。

% [Tutor Note: 原文 p015 第一条作「コンパクトかつ連結成分の時」，表述不整（按文意或为「連結」之衍字）；此处照原文译出，请审校。]

\section{热带仿射流形}

由适当的矩阵 $A\in\operatorname{Mat}(m,n;\R)$ 与向量 $b\in\R^n$ 表为
\begin{equation}
    f(x)=Ax+b \tag{9.1}
\end{equation}

% CHUNK c03 DONE pages p014–p015
